\documentclass[11pt,a4paper]{article}
\usepackage[utf8]{inputenc}
\usepackage[T1]{fontenc}
\usepackage[margin=1in]{geometry}
\usepackage{xcolor}
\usepackage{hyperref}
\definecolor{primarycolor}{RGB}{33,37,41}
\definecolor{secondarycolor}{RGB}{0,102,204}
\hypersetup{colorlinks=true,linkcolor=secondarycolor,urlcolor=secondarycolor,citecolor=secondarycolor}
\pagestyle{empty}
\setlength{\parindent}{0pt}
\setlength{\parskip}{0.85em}
\begin{document}
{\large\bfseries\color{primarycolor} Katarzyna Wiśniewska}\par
Sozialstation Dreisamtal gGmbH\par
\textbf{Bewerbung als Pflegehelferin für Ihren ambulanten Pflegedienst}\par
Sehr geehrte Damen und Herren,

Erfahrung in der stationären Altenpflege, ein LG1-Kurs und Deutschkenntnisse auf B2-Niveau sind meine Grundlagen für die Bewerbung bei Ihnen. Seit 2021 arbeite ich als Pflegehelferin im Seniorenzentrum Am Lindenhof. Zuvor war ich fünf Jahre als Krankenschwester auf einer Station für Innere Medizin in Polen tätig. Ihr Leitsatz „Pflege mit Zeit und Respekt“ passt zu meinem geduldigen Umgang mit älteren Menschen.

Im Seniorenzentrum leite ich seit 2022 wöchentlich eine Singgruppe für 6 bis 8 Menschen mit Demenz. Außerdem habe ich bisher 7 Pflegehelferinnen und Praktikanten nach dem Einarbeitungsplan der Einrichtung eingearbeitet. Bei einem Projekt zur Sturzprophylaxe wertete ich mit der Wohnbereichsleitung Sturzprotokolle aus und wir führten gezielte Maßnahmen ein. Im folgenden Halbjahr gingen die Stürze von 14 auf 8 zurück. Diese Erfahrungen möchte ich in Ihr Pflegeteam einbringen.

Meine bisherige Pflegepraxis liegt im stationären Bereich. Ihre zweiwöchige begleitete Einarbeitung bietet einen konkreten Rahmen für den Wechsel in die ambulante Pflege. Für die Pflegedokumentation nutze ich Vivendi PD sicher; mit MediFox habe ich noch nicht gearbeitet. Mein polnischer Pflegeabschluss befindet sich im Anerkennungsverfahren. Der Defizitbescheid liegt vor, die Kenntnisprüfung ist für März 2027 vorgesehen. Bis zur Anerkennung bewerbe ich mich ausdrücklich als Pflegehelferin.

Ich bevorzuge einen Beschäftigungsumfang von 75 \% und kann an jedem zweiten Wochenende arbeiten. Den vorausgesetzten Führerschein Klasse B habe ich noch nicht: Die Theorieprüfung ist bestanden, die praktische Prüfung ist für Januar 2027 vorgesehen. Derzeit nutze ich Fahrrad oder E-Bike. Gern möchte ich mit Ihnen besprechen, ob unter diesen Voraussetzungen eine Mitarbeit möglich ist. Über eine Einladung zu einem persönlichen Gespräch freue ich mich.

Mit freundlichen Grüßen\par
Katarzyna Wiśniewska
\end{document}